\sectionkicker{Source-backed technical notes}
\subsection*{振る舞いを教える――FLAN、T0、人間フィードバック: Technical Notes}
本欄は記事本文で圧縮した一次資料上の情報を、比較・再検証しやすい形で整理したものである。時系列、確認済みの主張、留意点、一次資料URLを掲載する。
\smallskip
\noindent{\small\textit{Technical Notesでは、一次資料から確認した公開・提供・研究上の事実を記載する。数値・能力評価や提供元・著者による比較表現は、独立再現済みの結果ではなく帰属claimとして扱う。}}\par
\medskip
\subsection*{Theme at a glance}
\addcontentsline{toc}{subsection}{Theme at a glance}
\begin{center}
\footnotesize
\begin{tabularx}{\linewidth}{@{}p{0.20\linewidth}p{0.10\linewidth}p{0.14\linewidth}X@{}}
\toprule
資料 & 位置づけ & 種別 & 時系列 \\
\midrule
FLAN & 主要資料 & 論文 & 2021-09-03 \\
T0 & 主要資料 & 論文 & 2021-10-15 \\
Recursive Book Summarization with Human Feedback & 主要資料 & 論文 & 2021-09-22, 2021-09-23 \\
Training Verifiers / GSM8K & 補足資料 & 論文 & 2021-10-27, 2021-10-29 \\
\bottomrule
\end{tabularx}
\end{center}
\normalsize

\begin{technicalnote}{FLAN}{主要資料}
\begingroup
% reader-facing Technical Notes break policy
\widowpenalty=10000
\clubpenalty=10000
\displaywidowpenalty=10000
\begin{tabularx}{\linewidth}{@{}>{\bfseries}p{0.22\linewidth}X@{}}
組織 & Google / authors \\
種別 & 論文 \\
時系列 & 2021-09-03 \\
\end{tabularx}
\smallskip
\begin{minipage}{\linewidth}
% reader-facing Technical Notes event-only fact/source tail group
{\bfseries 一次資料から整理したtechnical points}
\begin{itemize}[leftmargin=1.5em,itemsep=0.35em]
\item \textbf{一次情報で確認できる事実}: Google / authorsの選定済み一次資料では、FLANはpretrained language modelを、natural-language instruction templateで記述した60超のNLP task群に対してinstruction tuningし、未学習task typeへのzero-shot generalizationを評価した。 著者らは137B parameter modelを用い、finetuning dataset数・model scale・natural-language instructionの寄与をablationで調べた。
\end{itemize}
\begin{minipage}{\linewidth}
% reader-facing Technical Notes generic boundary/source tail group
\begin{samepage}
{\bfseries 一次資料}
\begingroup\sloppy
\Urlmuskip=0mu plus 2mu\relax
\begin{itemize}[leftmargin=1.5em,itemsep=0.25em]
\item {\scriptsize\url{http://arxiv.org/abs/2109.01652v5}}
\end{itemize}
\endgroup
\end{samepage}
\end{minipage}
\end{minipage}
\endgroup
\end{technicalnote}
\medskip

\begin{technicalnote}{T0}{主要資料}
\begingroup
% reader-facing Technical Notes break policy
\widowpenalty=10000
\clubpenalty=10000
\displaywidowpenalty=10000
\begin{tabularx}{\linewidth}{@{}>{\bfseries}p{0.22\linewidth}X@{}}
組織 & BigScience / authors \\
種別 & 論文 \\
時系列 & 2021-10-15 \\
\end{tabularx}
\smallskip
\begin{minipage}{\linewidth}
% reader-facing Technical Notes event-only fact/source tail group
{\bfseries 一次資料から整理したtechnical points}
\begin{itemize}[leftmargin=1.5em,itemsep=0.35em]
\item \textbf{一次情報で確認できる事実}: BigScience / authorsの選定済み一次資料では、T0は複数のsupervised datasetを人間可読なprompt形式へ変換し、datasetごとに複数の表現を持つmultitask mixtureでpretrained encoder-decoder modelをfine-tuneする構成を採った。 評価ではtraining mixtureから完全にheld-outしたtaskへのzero-shot generalizationを中心に、prompted multitask trainingの効果を比較した。
\end{itemize}
\begin{minipage}{\linewidth}
% reader-facing Technical Notes generic boundary/source tail group
\begin{samepage}
{\bfseries 一次資料}
\begingroup\sloppy
\Urlmuskip=0mu plus 2mu\relax
\begin{itemize}[leftmargin=1.5em,itemsep=0.25em]
\item {\scriptsize\url{http://arxiv.org/abs/2110.08207v3}}
\end{itemize}
\endgroup
\end{samepage}
\end{minipage}
\end{minipage}
\endgroup
\end{technicalnote}
\medskip

\begin{technicalnote}{Recursive Book Summarization with Human Feedback}{主要資料}
\begingroup
% reader-facing Technical Notes break policy
\widowpenalty=10000
\clubpenalty=10000
\displaywidowpenalty=10000
\begin{tabularx}{\linewidth}{@{}>{\bfseries}p{0.22\linewidth}X@{}}
組織 & OpenAI / authors \\
種別 & 論文 \\
時系列 & 2021-09-22, 2021-09-23 \\
\end{tabularx}
\smallskip
\begin{minipage}{\linewidth}
% reader-facing Technical Notes event-only fact/source tail group
{\bfseries 一次資料から整理したtechnical points}
\begin{itemize}[leftmargin=1.5em,itemsep=0.35em]
\item \textbf{一次情報で確認できる事実}: OpenAI / authorsの選定済み一次資料では、book-length summarizationを小区間へ分解し、小区間のsummaryを再帰的に要約するtask decompositionとhuman feedbackを組み合わせ、context windowを超える長文を扱う構成を検証した。 GPT-3をbehavioral cloningとreward modelingでfine-tuneし、人間labelerによるdemonstration/comparisonを使ってrecursive summarizationを学習させた。
\end{itemize}
\begin{minipage}{\linewidth}
% reader-facing Technical Notes generic boundary/source tail group
\begin{samepage}
{\bfseries 一次資料}
\begingroup\sloppy
\Urlmuskip=0mu plus 2mu\relax
\begin{itemize}[leftmargin=1.5em,itemsep=0.25em]
\item {\scriptsize\url{http://arxiv.org/abs/2109.10862v2}}
\item {\scriptsize\url{https://openai.com/index/summarizing-books}}
\end{itemize}
\endgroup
\end{samepage}
\end{minipage}
\end{minipage}
\endgroup
\end{technicalnote}
\medskip

\begin{technicalnote}{Training Verifiers / GSM8K}{補足資料}
\begingroup
% reader-facing Technical Notes break policy
\widowpenalty=10000
\clubpenalty=10000
\displaywidowpenalty=10000
\begin{tabularx}{\linewidth}{@{}>{\bfseries}p{0.22\linewidth}X@{}}
組織 & OpenAI / authors \\
種別 & 論文 \\
時系列 & 2021-10-27, 2021-10-29 \\
\end{tabularx}
\smallskip
\begin{minipage}{\linewidth}
% reader-facing Technical Notes event-only fact/source tail group
{\bfseries 一次資料から整理したtechnical points}
\begin{itemize}[leftmargin=1.5em,itemsep=0.35em]
\item \textbf{一次情報で確認できる事実}: OpenAI / authorsの選定済み一次資料では、GSM8Kは8.5K件のgrade-school math word problemを収録し、多段の数学推論を評価するdatasetとして提示された。 手法側ではmodelから複数のcandidate solutionを生成し、completionの正しさを判定するverifierでrankingして最上位を選ぶ構成を評価した。
\end{itemize}
\begin{minipage}{\linewidth}
% reader-facing Technical Notes generic boundary/source tail group
\begin{samepage}
{\bfseries 一次資料}
\begingroup\sloppy
\Urlmuskip=0mu plus 2mu\relax
\begin{itemize}[leftmargin=1.5em,itemsep=0.25em]
\item {\scriptsize\url{http://arxiv.org/abs/2110.14168v2}}
\item {\scriptsize\url{https://openai.com/index/solving-math-word-problems}}
\end{itemize}
\endgroup
\end{samepage}
\end{minipage}
\end{minipage}
\endgroup
\end{technicalnote}
\medskip
